\documentclass[9pt]{beamer}

% There are many different themes available for Beamer. A comprehensive
% list with examples is given here:
% http://deic.uab.es/~iblanes/beamer_gallery/index_by_theme.html

\usepackage[english]{babel}
\usepackage[TS1,T1]{fontenc}
\usepackage{array, booktabs}
\usepackage{multirow}
\usepackage[flushleft]{threeparttable}
\usepackage[font=small,labelfont=bf]{caption}
\usepackage[utf8]{inputenc}
\usepackage{bookman}
\usepackage{mathtools}
\usepackage{kotex}
\usepackage{adjustbox} % for \adjincludegraphics
\usepackage{textcomp}

\newcommand\ytl[2]{\parbox[b]{8em}{\hfill{\color{cyan}\bfseries\sffamily #1}~$\cdot$~}\makebox[0pt][c]{$\bullet$}\vrule\quad \parbox[c]{4.5cm}{\vspace{7pt}\color{red!40!black!80}\raggedright\sffamily #2.\\[7pt]}\\[-3pt]}

\usefonttheme{professionalfonts}

\usetheme{metropolis}

\setbeamertemplate{caption}[numbered]{}
\setbeamertemplate{navigation symbols}{}
\setbeamertemplate{footline}[frame number]

%\setbeameroption{show notes}

\graphicspath{{../../assets/img}}

\title{Structural Stability}

% A subtitle is optional and this may be deleted
%\subtitle{}

\author{Sung-Yong~Kim\inst{1}}
% - Give the names in the same order as the appear in the paper.
% - Use the \inst{?} command only if the authors have different
%   affiliation.

\institute[Changwon National University] % (optional, but mostly needed)
{
	\inst{1}%
	Associate Professor, School of Architecture\\
	Changwon National University}
% - Use the \inst command only if there are several affiliations.
% - Keep it simple, no one is interested in your street address.

\date{\today}
% - Either use conference name or its abbreviation.
% - Not really informative to the audience, more for people (including
%   yourself) who are reading the slides online

\subject{Lecture Notes}
% This is only inserted into the PDF information catalog. Can be left
% out. 

% If you have a file called "university-logo-filename.xxx", where xxx
% is a graphic format that can be processed by latex or pdflatex,
% resp., then you can add a logo as follows:

% \pgfdeclareimage[height=0.5cm]{university-logo}{university-logo-filename}
% \logo{\pgfuseimage{university-logo}}


% Let's get started
\begin{document}

\begin{frame}
	\titlepage
\end{frame}

\section{AISC-LRFD Design Curve}

\begin{frame}{AISC-LRFD Design Curve}
	Euler Buckling Load:
\[P_\mathrm{cr} = \frac{\pi^2 EI}{(KL)^2}\]
	Buckling Strength (탄성좌굴강도):
\[F_\mathrm{cr} = \frac{\pi^2 EI}{A(KL)^2} = \frac{\pi^2 E}{(KL/r)^2}\]
where $r$: 단면2차반경
\[r = \sqrt{I/A}\]
	항복강도($F_y$)와 탄성좌굴강도($F_\mathrm{cr}$)의 교차점: $F_y=F_\mathrm{cr}$에 대한 $KL/r$ 값
\[F_\mathrm{cr} = F_y \Rightarrow \frac{\pi^2 E}{(KL/r)^2} = F_y \Rightarrow (KL/r)_i = \pi\sqrt{\frac{E}{F_y}}\]
\end{frame}

\begin{frame}{AISC-LRFD Design Curve and KDS Formulation}
	Dimensionless Slenderness Ratio (무차원세장비):
\[\lambda_c = \frac{KL}{r}\sqrt{\frac{F_y}{\pi^2 E}}\]
	Design Curve:
\[\frac{F}{F_y} = \begin{cases}\mathrm{exp}(-0.419\lambda_c^2) & \lambda_c \leq 1.5 \\
							   0.877\lambda_c^{-2} & \lambda_c>1.5\end{cases}\]
	KDS Formulation:
\[\lambda_c = \frac{KL}{r}\sqrt{\frac{F_y}{\pi^2 E}} = \sqrt{\frac{F_y}{F_e}} \leq 1.5 \Rightarrow \]
\end{frame}
\end{document}
